\section{Introduction}

% 1--2 short paragraphs:
% - The RNA--RNA interaction prediction problem; why fast, accurate scanning matters
%   (siRNA design, ncRNA target screens).
% - RIsearch2 baseline (cite); identified pain points: install friction, monolithic CLI,
%   no Python integration, off-target pipeline written against an older Python.
% - One sentence stating the contribution: RIsearch3 — Rust core + Python bindings +
%   modernised pipeline.

\section{Implementation}

\subsection{Rust core}\label{subsec:core}

RIsearch3 reimplements the suffix-array seed-and-extend algorithm of
RIsearch2~\citep{alkan2017risearch2} in Rust. The two-stage scheme is
preserved without modification: a suffix array of the concatenated target
sequences (with the reverse-complement strand stored explicitly) is
traversed in parallel with a per-query suffix array to locate maximal seed
matches, and each seed is then extended by dynamic programming under the
simplified nearest-neighbour energy model of~\citet{wenzel2012risearch}.
Match identity follows the parallel-direction trick of
RIsearch2~\citep[Fig.~1B]{alkan2017risearch2}, with $G{-}A$ and $U{-}C$ identities
counted as matches so that $G{-}U$ wobble pairs are admitted.

The rewrite focuses on engineering rather than algorithmic change. Suffix-array
construction is delegated to the \texttt{libsais} library, replacing the
\texttt{libdivsufsort} library used by RIsearch2; an optional OpenMP build
provides parallel index construction. Search is parallelised with Rayon
across seed groups, with each worker holding its own dynamic-programming
context to avoid contention. The DP extension uses an X-drop early-termination
heuristic, evaluates the energy model inline rather than through precomputed
lookup tables, and represents the forward and reverse strands through Rust
const generics so that the compiler specialises both code paths separately.
The crate uses the jemalloc allocator. Both energy parameter sets supported
by RIsearch2 (\texttt{t04} and \texttt{t99}, derived from the Turner
nearest-neighbour parameters) are retained.

Behaviour preservation is enforced by a parity test suite. Integration tests
compare RIsearch3 output against the original C implementation across seed
length and position constraints, mismatch settings, extension parameters,
both energy parameter sets, $G{-}U$ wobble handling, and the optional banding
mode. Outputs match bit-for-bit on the regression corpus. Migration risk is
therefore minimal: existing analyses can adopt RIsearch3 without
re-validating downstream results.

On equivalent inputs, RIsearch3 is $2{-}3\times$ faster than RIsearch2 in
wall-clock time.% TODO: name the benchmark dataset (size, hardware) once finalised.

\subsection{Command-line interface}\label{subsec:cli}

% - Subcommand structure, sensible defaults, useful --help.
% - One short usage example in a verbatim block.

\subsection{Python bindings}\label{subsec:bindings}

% - PyO3-based bindings; pip-installable wheel.
% - Minimal Python snippet showing a query + result iteration.
% - Sets up the off-target pipeline as the worked example.

\subsection{siRNA off-target pipeline}\label{subsec:pipeline}

% - Rewritten in modern Python; Polars for tabular ops.
% - YAML configuration; single CLI entry point.
% - Built on top of the Python bindings (Sec.~\ref{subsec:bindings}) — concrete
%   demonstration of how downstream tools can consume RIsearch3.
% - Inline benchmark: end-to-end runtime on a reference siRNA set vs. the
%   RIsearch2-era pipeline. Reference Fig.~\ref{fig:pipeline_benchmark}.

\section{Conclusion}

% 1 short paragraph: what's available, who benefits, what's next (additional energy
% models, further integration targets).
